\chapter{Limitations}
\label{ch:limitations}

While the MEDF platform represents a significant step towards operationalizing AI ethics, it is essential to acknowledge its limitations. These limitations define the boundaries of the current implementation and highlight important areas for future research and development. The key limitations can be grouped into three categories: subjectivity of inputs, simplification of models, and scope of evaluation.

\section{Subjectivity of Inputs}
\label{sec:subjectivity}

The adage \enquote{garbage in, garbage out} applies forcefully to the MEDF framework. The validity of the platform's output is fundamentally dependent on the quality and integrity of its inputs.

\begin{itemize}[leftmargin=2em]
  \item \textbf{Baseline Score Subjectivity.} The initial dimension scores for an AI system are a critical input. In the current implementation, these scores are assumed to be provided by the user (e.g., the development team). These scores are inherently subjective and can be influenced by the assessor's biases, knowledge, and incentives. A team might be motivated to provide inflated scores to make their system appear more ethical than it is. While MEDF provides a framework for analyzing these scores, it does not have a built-in mechanism for independently verifying their accuracy against ground truth.

  \item \textbf{Stakeholder Weight Subjectivity.} Similarly, the stakeholder weights are also subjective. While they are intended to represent the genuine priorities of different groups, determining the \enquote{correct} weights for a group like the \enquote{Affected Community} is a significant challenge in itself. The default profiles provided in MEDF are based on reasonable assumptions, but they are not the result of empirical sociological research. A comprehensive and truly representative evaluation would require a more rigorous process for eliciting these weights from the actual stakeholder groups.
\end{itemize}

\section{Simplification of Models}
\label{sec:simplification}

To create a tractable computational framework, MEDF necessarily simplifies the complex reality of ethical decision-making.

\begin{itemize}[leftmargin=2em]
  \item \textbf{Linearity and Independence Assumptions.} The scoring models (TOPSIS and WSM) and the weighting schemes implicitly assume a degree of independence between the ethical dimensions. In reality, these dimensions are often deeply intertwined. For example, improving transparency might have a positive impact on accountability but could have a negative impact on safety if it exposes system vulnerabilities. The current models do not explicitly capture these complex, non-linear interactions.

  \item \textbf{Reduction of Ethics to a Score.} The act of reducing a complex ethical profile to a single numerical score is a powerful simplification, but it is also a lossy one. A single score cannot capture the full nuance of an ethical dilemma. Two systems could receive the same overall score but have vastly different ethical profiles: one might be moderately weak across all dimensions, while the other might excel in five dimensions but have a critical flaw in one. While MEDF provides dimensional breakdowns to mitigate this, the focus on a single score can risk oversimplification.
\end{itemize}

\section{Scope of Evaluation}
\label{sec:scope}

The current implementation of MEDF has a defined scope, and several important aspects of AI governance are not yet included.

\begin{itemize}[leftmargin=2em]
  \item \textbf{No Primary Research.} The project mandate explicitly excluded primary research such as user surveys, interviews, or statistical user studies. As a result, the framework's usability and the validity of its default assumptions have not been empirically tested with real-world users.

  \item \textbf{Focus on Post-Hoc Evaluation.} MEDF is primarily designed for the evaluation of existing systems or well-defined designs. While it can be used to inform the design process, it is not a generative tool that actively assists in the creation of more ethical designs from scratch. It is an evaluative, not a creative, framework.

  \item \textbf{Limited Harm Taxonomy.} The harm assessment module provides a high-level estimation of risk. A more sophisticated implementation would incorporate a more detailed and evidence-based taxonomy of AI-related harms, potentially drawing on fields like law and social science to create a more granular risk model.

  \item \textbf{Incomplete Separation of Concerns.} The conflict analysis logic is currently implemented within the API router (\texttt{routers/conflicts.py}) rather than an extracted standalone engine module. While the functionality is complete and all tests pass, this architectural shortcut represents technical debt. A future refactoring should separate the conflict analysis logic into a dedicated module to complete the intended separation of concerns.
\end{itemize}

These limitations do not invalidate the MEDF approach. Rather, they clarify its role as a specific type of decision-support tool and provide a clear roadmap for future work aimed at building more comprehensive, empirically-grounded, and nuanced tools for computational ethics.
